\scanpage{19}
- \mspage{15} -
\[
(6.7)\qquad y^2=x^3-x\qquad(\text{si }c_4=c_6=0)
\]
c) $p=2$ : ici, on met la courbe sour la forme (6.3) et on résoud (6.4) avec $a_3=1$, $a_4=a_6=0$ pour obtenir la forme type
\[
(6.8)\qquad y^2+y=x^3\qquad(\text{si }c_4=c_6=0).
\]
Pour compléter le tableau, reste à décrire les automorphismes des courbes elliptiques.

\uline{Proposition} 6.9.\quad\uline{Soit} $E$ \uline{une courbe de genre un sur une base} $S$, \uline{sans fibre additive}

(I) \uline{Si} $c_4$ et $c_6$ \uline{sont inversibles, les seuls automorphismes de} $E$ \uline{sont l'identité et la symétrie}.
\margin{[vrai seulement en faisceautisant !]}

(II) \uline{Si} 2 et 3 \uline{sont inversibles, l'application canonique}
\[
(6.6.)\qquad\diagram{\begin{tikzcd}\uline{\mathrm{Aut}}(E)\arrow[r]&G_m\end{tikzcd}}
\]
\uline{donnée par l'action sur les différentielles invariantes, est une immersion fermée}.

(III) \uline{Si} 2 \uline{et} 3 \uline{sont inversibles}, \uline{et} $c_4=0$ (resp. $c_6=0$) \uline{l'image de} (6.6) \uline{est} $\mu_6$ (resp. $\mu_4$).

\uline{Supposons maintenant que} $S$ \uline{soit le spectre d'un corps}

(IV) \uline{Si} $3=c_4=c_6=0$, $\mathrm{Aut}(E)$ \uline{est géométriquement \unsure{isomorphe} au produit semi-direct} (\uline{non commutatif}) \uline{de} $\mathbb Z/(4)$ \uline{par} $\mathbb Z/(3)$, \uline{et l'image de} (6.6) \uline{est} $\mu_4$.

(V) \uline{Si} $2=c_4=c_6=0$, $\mathrm{Aut}(E)$ \uline{est géométriquement d'ordre} 24, \uline{et l'image de} (6.6) \uline{est} $\mu_3$.
\margin{\unsure{$\rho$}}
\margin{18}
